% VI37-DETAILED-PROBLEM-16
\begin{problem}[VI.37.P16: Two homogeneous coordinate rings for \(\mathbb P^1\)]\label{prob:vi37-16}
Prove
\[
\Proj k[s,t]
\cong
\Proj k[X,Y,Z]/(XZ-Y^2)
\]
using the second Veronese ring.

\textbf{Provenance.} Canonical VI/37 extension.
\end{problem}

\ifdefined\FullProblemDossiers
\begin{solution}
Let
\[
R=k[s,t].
\]
Its second Veronese ring is
\[
R^{(2)}
=
\bigoplus_{m\ge0}R_{2m}.
\]
It is generated in its regraded degree one by
\[
s^2,\quad st,\quad t^2.
\]
Thus the map
\[
k[X,Y,Z]\longrightarrow R^{(2)},
\qquad
X\mapsto s^2,\quad Y\mapsto st,\quad Z\mapsto t^2
\]
is surjective. The relation
\[
XZ-Y^2
\]
lies in the kernel, and the standard monomial argument shows it is the only defining relation:
\[
R^{(2)}
\cong
k[X,Y,Z]/(XZ-Y^2).
\]

VI/35 proves Veronese invariance:
\[
\Proj R\cong\Proj R^{(2)}.
\]
Therefore
\[
\boxed{
\Proj k[s,t]
\cong
\Proj k[X,Y,Z]/(XZ-Y^2).
}
\]
This exhibits explicitly how the same abstract projective scheme can carry different homogeneous coordinate rings arising from different embeddings.
\end{solution}
\fi
